\begin{abstract}
Software development organizations generate large volumes of process data through \textit{issue} management systems, version control repositories, continuous integration \textit{pipelines}, and code review platforms. Despite this extensive instrumentation, the prediction of software development process behavior still relies predominantly on aggregated metrics and heuristic estimation methods that treat development workflows as black boxes. As a result, existing approaches provide a limited understanding of the structural dynamics and path-dependent variability that characterize real software development processes. This thesis investigates whether models derived from process mining can provide a more informative basis for the probabilistic forecasting of software development processes. In particular, the research tests the hypothesis that process-derived features and explicit stochastic models can improve predictive performance in comparison with approaches based exclusively on traditional software repository metrics. To address this problem, the thesis proposes PATHCAST (PATH-aware foreCASTing), a computational method that integrates process mining, Markov chain modeling, Monte Carlo simulation, and machine learning into a unified analytical \textit{pipeline}. The approach transforms heterogeneous data from software repositories into structured event logs, from which process models are discovered to represent actual development flows. These models are then used to derive stochastic state-transition models that enable the probabilistic simulation of process execution trajectories and their outcome distributions. The research follows the \textit{Design Science Research} methodology, in which the proposed method is implemented and empirically evaluated based on historical data from software projects. The evaluation compares the predictive performance of PATHCAST with baseline models based on repository metrics and conventional forecasting techniques. The expected results are that the explicit incorporation of process structure into forecasting models will improve both predictive accuracy and interpretability, thereby contributing to the advancement of process intelligence applied to software engineering.
	\keywords{Process Mining, Probabilistic Forecasting, Markov Chains, Monte Carlo Simulation, Machine Learning, Design Science Research}
\end{abstract}
